\documentclass[10pt,a4paper]{amsart}
\usepackage{amsmath,amssymb,amsthm,graphicx}
\usepackage{verbatim}
\newenvironment{reference}{\comment}{\endcomment}
%\newenvironment{reference}{}{}
\newenvironment{slogan}{\comment}{\endcomment}
\newenvironment{history}{\comment}{\endcomment}

% For commutative diagrams we use Xy-pic
\usepackage[all]{xy}

% We use 2cell for 2-commutative diagrams.
\xyoption{2cell}
% Theorem environments.
%
\theoremstyle{plain}
\newtheorem{theorem}[subsection]{Theorem}
\newtheorem{proposition}[subsection]{Proposition}
\newtheorem{lemma}[subsection]{Lemma}

\theoremstyle{definition}
\newtheorem{definition}[subsection]{Definition}
\newtheorem{example}[subsection]{Example}
\newtheorem{exercise}[subsection]{Exercise}
\newtheorem{situation}[subsection]{Situation}

\theoremstyle{remark}
\newtheorem{remark}[subsection]{Remark}
\newtheorem{remarks}[subsection]{Remarks}

\numberwithin{equation}{subsection}

% Macros
%
\def\lim{\mathop{\mathrm{lim}}\nolimits}
\def\colim{\mathop{\mathrm{colim}}\nolimits}
\def\Spec{\mathop{\mathrm{Spec}}}
\def\Hom{\mathop{\mathrm{Hom}}\nolimits}
\def\Ext{\mathop{\mathrm{Ext}}\nolimits}
\def\SheafHom{\mathop{\mathcal{H}\!\mathit{om}}\nolimits}
\def\SheafExt{\mathop{\mathcal{E}\!\mathit{xt}}\nolimits}
\def\Sch{\mathit{Sch}}
\def\Mor{\mathop{\mathrm{Mor}}\nolimits}
\def\Ob{\mathop{\mathrm{Ob}}\nolimits}
\def\Sh{\mathop{\mathit{Sh}}\nolimits}
\def\NL{\mathop{N\!L}\nolimits}
\def\CH{\mathop{\mathrm{CH}}\nolimits}
\def\proetale{{pro\text{-}\acute{e}tale}}
\def\etale{{\acute{e}tale}}
\def\QCoh{\mathit{QCoh}}
\def\Ker{\mathop{\mathrm{Ker}}}
\def\Im{\mathop{\mathrm{Im}}}
\def\Coker{\mathop{\mathrm{Coker}}}
\def\Coim{\mathop{\mathrm{Coim}}}

% Boxtimes
%
\DeclareMathSymbol{\boxtimes}{\mathbin}{AMSa}{"02}

%
% Macros for moduli stacks/spaces
%
\def\QCohstack{\mathcal{QC}\!\mathit{oh}}
\def\Cohstack{\mathcal{C}\!\mathit{oh}}
\def\Spacesstack{\mathcal{S}\!\mathit{paces}}
\def\Quotfunctor{\mathrm{Quot}}
\def\Hilbfunctor{\mathrm{Hilb}}
\def\Curvesstack{\mathcal{C}\!\mathit{urves}}
\def\Polarizedstack{\mathcal{P}\!\mathit{olarized}}
\def\Complexesstack{\mathcal{C}\!\mathit{omplexes}}
% \Pic is the operator that assigns to X its picard group, usage \Pic(X)
% \Picardstack_{X/B} denotes the Picard stack of X over B
% \Picardfunctor_{X/B} denotes the Picard functor of X over B
\def\Pic{\mathop{\mathrm{Pic}}\nolimits}
\def\Picardstack{\mathcal{P}\!\mathit{ic}}
\def\Picardfunctor{\mathrm{Pic}}
\def\Deformationcategory{\mathcal{D}\!\mathit{ef}}
\usepackage{cleveref}
\crefname{originallemma}{original lemma}{original lemmas}
\Crefname{originallemma}{Original lemma}{Original lemmas}
\makeatletter
\newlabel{divisors-lemma-reflexive-over-regular-dim-2}{{0B3N}{source}{}{sourcepointer-1}{}}
\newlabel{divisors-lemma-reflexive-over-regular-dim-2@cref}{{[originallemma][0][]0B3N}{[1][0][]source}{}{}{}}
\makeatother
\crefname{originallemma}{original lemma}{original lemmas}
\Crefname{originallemma}{Original lemma}{Original lemmas}
\makeatletter
\newlabel{divisors-lemma-wedge-product-ses}{{0B3T}{source}{}{sourcepointer-2}{}}
\newlabel{divisors-lemma-wedge-product-ses@cref}{{[originallemma][0][]0B3T}{[1][0][]source}{}{}{}}
\makeatother
\crefname{originalsection}{original section}{original sections}
\Crefname{originalsection}{Original section}{Original sections}
\makeatletter
\newlabel{duality-section-duality}{{0AU3}{source}{}{sourcepointer-3}{}}
\newlabel{duality-section-duality@cref}{{[originalsection][0][]0AU3}{[1][0][]source}{}{}{}}
\makeatother
\crefname{originallemma}{original lemma}{original lemmas}
\Crefname{originallemma}{Original lemma}{Original lemmas}
\makeatletter
\newlabel{dualizing-lemma-regular-gorenstein}{{0AWX}{source}{}{sourcepointer-4}{}}
\newlabel{dualizing-lemma-regular-gorenstein@cref}{{[originallemma][0][]0AWX}{[1][0][]source}{}{}{}}
\makeatother
\crefname{originalreference}{original reference}{original references}
\Crefname{originalreference}{Original reference}{Original references}
\makeatletter
\newlabel{lemma-blowup-regular}{{0AGR}{source}{}{sourcepointer-5}{}}
\newlabel{lemma-blowup-regular@cref}{{[originalreference][0][]0AGR}{[1][0][]source}{}{}{}}
\makeatother
\crefname{originalreference}{original reference}{original references}
\Crefname{originalreference}{Original reference}{Original references}
\makeatletter
\newlabel{lemma-cohomology-of-blowup}{{0AGS}{source}{}{sourcepointer-6}{}}
\newlabel{lemma-cohomology-of-blowup@cref}{{[originalreference][0][]0AGS}{[1][0][]source}{}{}{}}
\makeatother
\crefname{originalreference}{original reference}{original references}
\Crefname{originalreference}{Original reference}{Original references}
\makeatletter
\newlabel{lemma-differentials-of-blowup}{{0B4L}{source}{}{sourcepointer-7}{}}
\newlabel{lemma-differentials-of-blowup@cref}{{[originalreference][0][]0B4L}{[1][0][]source}{}{}{}}
\makeatother
\crefname{originalreference}{original reference}{original references}
\Crefname{originalreference}{Original reference}{Original references}
\makeatletter
\newlabel{remark-dualizing-setup}{{0B4R}{source}{}{sourcepointer-8}{}}
\newlabel{remark-dualizing-setup@cref}{{[originalreference][0][]0B4R}{[1][0][]source}{}{}{}}
\makeatother
\title{Comparing top relative differentials with the dualizing module along sequences of point blowups}
\author{The Stacks Project Authors}
\date{stacks-project:resolve.tex:local-snapshot-20261012; GNU FDL 1.2 or later}
\begin{document}
\maketitle
\noindent\textit{Editorial note.} All delivered mathematical text is the author's own, copied from the pinned Stacks Project chapter file resolve.tex; no statement or proof step was rewritten, reordered or removed.\par\medskip
\noindent\textit{Editorial note.} The counted proof is the author's complete argument: the duality reduction of the required map to a map of derived pushforwards, the local computation identifying the top relative differentials of a point blowup up to a twist by the exceptional divisor in the two possible cases, and the application of the cohomology computation for a point blowup to complete the induction.\par\medskip
\noindent\textit{Editorial note.} The label lemma-compare-differentials-dualizing names the counted statement, and the proof environment immediately following it in the source is delivered verbatim as its proof.\par\medskip
\noindent\textit{Editorial note.} The counted statement is carried by the author's own lemma environment in the chapter source.\par\medskip
\noindent\textit{Editorial note.} Pre-open gate: python operations/agent-authoring/tools/precard\_gate.py --source resolve.tex --label lemma-compare-differentials-dualizing --start 2290 --end 2346 --local-id 28095 returned RESULT=OK (L28095 free, occupancy FREE, public interval clear). Pre-assembly tag self-check: canonical tag resolved and not present in reference/public-items.jsonl.\par\medskip
\noindent\textit{Editorial note.} The statement and proof are copied from the pinned chapter file; the cited inputs are kept as external pointers with their canonical Stacks tags.\par\medskip
\noindent\textit{Allowed standard prerequisites (editorial).}\begin{itemize}
\item Regular rings are Gorenstein and blowups of regular surfaces are regular. Dualizing Complexes, the lemma that a regular ring is Gorenstein, and Same chapter, the lemma that a blowup of a regular surface in a closed point is regular. Supply the invertibility of the modules serving as dualizing modules for the base and for each term of the sequence.
\item Duality for a blowup and the identification of the dualizing module as an inverse image. Duality for Schemes, the section on duality and the discussion of the dualizing setup of the chapter. Supplies the translation of a map on the top relative differentials of the new term into a map of derived pushforwards, and the identification of the dualizing module of the new term as the inverse image of the previous one.
\item Local computation of the relative differentials of a point blowup and the wedge of a short exact sequence. Same chapter, the lemma computing the relative differentials of a blowup in a closed point, and Divisors, the lemma describing the wedge product of a short exact sequence of locally free sheaves. Supply the two possibilities for the top relative differentials of the new term, differing by a twist by the exceptional divisor.
\item Vanishing of the first pushforward of the twisting sheaves of a point blowup. Same chapter, the lemma computing the cohomology of the structure sheaf and of its twists on a point blowup. Supplies the vanishing of the first pushforward and the identification of the zeroth pushforward with the previous term.
\end{itemize}\medskip
\ifdefined\hypertarget\hypertarget{sourcepointer-1}{}\fi
\noindent\textit{Source reference.} Original lemma 0B3N: sources/stacks-chow/tags.txt, source lines 13971--13971. Prerequisite pointer invoked at proof line(s) 2327; it names a standard fact of a sibling chapter that the counted argument uses, and every step of the counted proof built on it is retained verbatim in the delivered body.\par\medskip
\ifdefined\hypertarget\hypertarget{sourcepointer-2}{}\fi
\noindent\textit{Source reference.} Original lemma 0B3T: sources/stacks-chow/tags.txt, source lines 13976--13976. Prerequisite pointer invoked at proof line(s) 2341; it names a standard fact of a sibling chapter that the counted argument uses, and every step of the counted proof built on it is retained verbatim in the delivered body.\par\medskip
\ifdefined\hypertarget\hypertarget{sourcepointer-3}{}\fi
\noindent\textit{Source reference.} Original section 0AU3: sources/stacks-chow/tags.txt, source lines 13636--13636. Prerequisite pointer invoked at proof line(s) 2321; it names a standard fact of a sibling chapter that the counted argument uses, and every step of the counted proof built on it is retained verbatim in the delivered body.\par\medskip
\ifdefined\hypertarget\hypertarget{sourcepointer-4}{}\fi
\noindent\textit{Source reference.} Original lemma 0AWX: sources/stacks-chow/tags.txt, source lines 13735--13735. Prerequisite pointer invoked at proof line(s) 2306; it names a standard fact of a sibling chapter that the counted argument uses, and every step of the counted proof built on it is retained verbatim in the delivered body.\par\medskip
\ifdefined\hypertarget\hypertarget{sourcepointer-5}{}\fi
\noindent\textit{Source reference.} Original reference 0AGR: sources/stacks-resolve/pinned-original-source.tex, source lines 727--727. Prerequisite pointer invoked at proof line(s) 2312; it names a standard fact of the same chapter that the counted argument uses, and every step of the counted proof built on it is retained verbatim in the delivered body.\par\medskip
\ifdefined\hypertarget\hypertarget{sourcepointer-6}{}\fi
\noindent\textit{Source reference.} Original reference 0AGS: sources/stacks-resolve/pinned-original-source.tex, source lines 786--786. Prerequisite pointer invoked at proof line(s) 2345; it names a standard fact of the same chapter that the counted argument uses, and every step of the counted proof built on it is retained verbatim in the delivered body.\par\medskip
\ifdefined\hypertarget\hypertarget{sourcepointer-7}{}\fi
\noindent\textit{Source reference.} Original reference 0B4L: sources/stacks-resolve/pinned-original-source.tex, source lines 1012--1012. Prerequisite pointer invoked at proof line(s) 2334; it names a standard fact of the same chapter that the counted argument uses, and every step of the counted proof built on it is retained verbatim in the delivered body.\par\medskip
\ifdefined\hypertarget\hypertarget{sourcepointer-8}{}\fi
\noindent\textit{Source reference.} Original reference 0B4R: sources/stacks-resolve/pinned-original-source.tex, source lines 1807--1807. Prerequisite pointer invoked at proof line(s) 2311; it names a standard fact of the same chapter that the counted argument uses, and every step of the counted proof built on it is retained verbatim in the delivered body.\par\medskip
\begin{lemma}
\label{lemma-compare-differentials-dualizing}
Let $p$ be a prime number.
Let $A$ be a regular local ring of dimension $2$ and characteristic $p$.
Let $A_0 \subset A$ be a subring such that $\Omega_{A/A_0}$ is free
of rank $r < \infty$. Set $\omega_A = \Omega^r_{A/A_0}$. If $X \to \Spec(A)$
is the result of a sequence of blowups in closed points, then
there exists a map
$$
\varphi_X : (\Omega^r_{X/\Spec(A_0)})^{**} \longrightarrow \omega_X
$$
extending the given identification in the generic point.
\end{lemma}

\begin{proof}
Observe that $A$ is Gorenstein (Dualizing Complexes,
Lemma \ref{dualizing-lemma-regular-gorenstein})
and hence the invertible module $\omega_A$ does indeed serve
as a dualizing module. Moreover, any $X$ as in the lemma
has an invertible dualizing module $\omega_X$ as $X$ is regular
(hence Gorenstein) and proper over $A$, see
Remark \ref{remark-dualizing-setup} and
Lemma \ref{lemma-blowup-regular}.
Suppose we have constructed the map
$\varphi_X : (\Omega^r_{X/A_0})^{**} \to \omega_X$
and suppose that $b : X' \to X$ is a blowup in a closed point.
Set $\Omega^r_X = (\Omega^r_{X/A_0})^{**}$ and
$\Omega^r_{X'} = (\Omega^r_{X'/A_0})^{**}$. Since $\omega_{X'} = b^!(\omega_X)$
a map $\Omega^r_{X'} \to \omega_{X'}$ is the same thing as a map
$Rb_*(\Omega^r_{X'}) \to \omega_X$. See discussion in
Remark \ref{remark-dualizing-setup} and
Duality for Schemes, Section \ref{duality-section-duality}.
Thus in turn it suffices to produce a map
$$
Rb_*(\Omega^r_{X'}) \longrightarrow \Omega^r_X
$$
The sheaves $\Omega^r_{X'}$ and $\Omega^r_X$ are invertible, see
Divisors, Lemma \ref{divisors-lemma-reflexive-over-regular-dim-2}.
Consider the exact sequence
$$
b^*\Omega_{X/A_0} \to \Omega_{X'/A_0} \to \Omega_{X'/X} \to 0
$$
A local calculation shows that $\Omega_{X'/X}$ is isomorphic
to an invertible module on the exceptional divisor $E$, see
Lemma \ref{lemma-differentials-of-blowup}. It follows that
either
$$
\Omega^r_{X'} \cong (b^*\Omega^r_X)(E)
\quad\text{or}\quad
\Omega^r_{X'} \cong b^*\Omega^r_X
$$
see Divisors, Lemma \ref{divisors-lemma-wedge-product-ses}.
(The second possibility never happens in characteristic zero, but
can happen in characteristic $p$.) In both cases we see that
$R^1b_*(\Omega^r_{X'}) = 0$ and $b_*(\Omega^r_{X'}) = \Omega^r_X$ by
Lemma \ref{lemma-cohomology-of-blowup}.
\end{proof}

\end{document}
